% !TEX program = xelatex
% =====================================================================
%  BÁO CÁO DỰ ÁN: FRACTIONAL ROYALTY NFT MARKETPLACE
%  Compiler: xelatex (khuyến nghị) hoặc pdflatex
%  Trên Overleaf: Menu → Settings → Compiler → XeLaTeX
% =====================================================================
\documentclass[12pt,a4paper,oneside]{report}

% ---------- Encoding & Vietnamese ----------
\usepackage{iftex}
\ifPDFTeX
  \usepackage[utf8]{inputenc}
  \usepackage[T5]{fontenc}
  \usepackage[vietnamese]{babel}
  \usepackage{lmodern}
\else
  \usepackage{fontspec}
  \usepackage{polyglossia}
  \setmainlanguage{vietnamese}
  \setmainfont{Latin Modern Roman}
  \setsansfont{Latin Modern Sans}
  \setmonofont{Latin Modern Mono}[Scale=0.88]
\fi

% ---------- Layout ----------
\usepackage[a4paper,top=2.5cm,bottom=2.5cm,left=2.8cm,right=2.5cm]{geometry}
\usepackage{setspace}\onehalfspacing
\usepackage{parskip}
\usepackage{microtype}
\usepackage{enumitem}
\setlist{nosep,leftmargin=1.5em}

% ---------- Math ----------
\usepackage{amsmath,amssymb,amsthm}

% ---------- Graphics, tables ----------
\usepackage{graphicx,float,xcolor}
\usepackage{booktabs,longtable,array,tabularx,multirow}
\usepackage{caption}\captionsetup{font=small,labelfont=bf}

% ---------- TikZ ----------
\usepackage{tikz}
\usetikzlibrary{arrows.meta,positioning,shapes.geometric,calc,decorations.pathreplacing}

% ---------- Code listings ----------
\usepackage{listings}
\definecolor{codebg}{RGB}{248,248,250}
\definecolor{codekw}{RGB}{60,60,180}
\definecolor{codestr}{RGB}{0,120,60}
\definecolor{codecom}{RGB}{128,128,128}
\definecolor{codetype}{RGB}{150,60,60}

\lstdefinelanguage{Solidity}{
  keywords={pragma,solidity,contract,interface,library,is,public,private,internal,external,
    pure,view,payable,returns,return,memory,storage,calldata,function,constructor,
    modifier,event,emit,indexed,mapping,struct,enum,require,revert,assert,
    if,else,for,while,new,delete,using,override,virtual,immutable,constant,
    true,false,this,unchecked,try,catch},
  keywordstyle=\color{codekw}\bfseries,
  ndkeywords={address,bool,string,bytes,bytes32,uint,uint96,uint256,int256},
  ndkeywordstyle=\color{codetype},
  sensitive=true,
  comment=[l]{//},morecomment=[s]{/*}{*/},
  commentstyle=\color{codecom}\itshape,
  stringstyle=\color{codestr},
  morestring=[b]",morestring=[b]'
}

\lstdefinelanguage{JavaScript}{
  keywords={async,await,const,let,var,function,return,if,else,for,new,this,
    class,typeof,try,catch,throw,true,false,null,undefined},
  keywordstyle=\color{codekw}\bfseries,
  sensitive=true,
  comment=[l]{//},morecomment=[s]{/*}{*/},
  commentstyle=\color{codecom}\itshape,
  stringstyle=\color{codestr},
  morestring=[b]",morestring=[b]',morestring=[b]`
}

\lstset{
  backgroundcolor=\color{codebg},
  basicstyle=\ttfamily\footnotesize,
  breaklines=true,columns=fullflexible,
  frame=single,framerule=0.3pt,rulecolor=\color{gray!50},
  numbers=left,numberstyle=\tiny\color{gray},numbersep=7pt,
  showstringspaces=false,tabsize=2,captionpos=b,
  xleftmargin=1em,escapeinside={(*@}{@*)},
}

% ---------- Theorems ----------
\theoremstyle{plain}
\newtheorem{theorem}{Định lý}[chapter]
\newtheorem{lemma}[theorem]{Bổ đề}
\newtheorem{proposition}[theorem]{Mệnh đề}
\theoremstyle{definition}
\newtheorem{definition}[theorem]{Định nghĩa}
\theoremstyle{remark}
\newtheorem{remark}[theorem]{Nhận xét}

% ---------- Hyperlinks ----------
\usepackage[hidelinks,colorlinks=true,linkcolor=blue!60!black,citecolor=teal!60!black,urlcolor=blue!60!black]{hyperref}

% ---------- Headers ----------
\usepackage{fancyhdr}
\pagestyle{fancy}\fancyhf{}
\renewcommand{\headrulewidth}{0.4pt}
\fancyhead[L]{\small\leftmark}
\fancyhead[R]{\small\thepage}

% =====================================================================
\begin{document}

% ==================== TRANG BÌA ====================
\begin{titlepage}
\centering
\vspace*{1.5cm}
{\Huge\bfseries Fractional Royalty\\[8pt] NFT Marketplace\par}
\vspace{0.8cm}
{\Large Báo cáo kỹ thuật dự án\par}
\vspace{0.4cm}
{\large\itshape ERC-721 $\bullet$ EIP-2981 Multi-Receiver $\bullet$ Pull-Payment $\bullet$ Atomic Swap\par}
\vspace{2cm}

\begin{tikzpicture}[node distance=1.6cm and 2.5cm,
  ct/.style={rectangle,rounded corners=6pt,draw,thick,minimum width=3.4cm,minimum height=1.1cm,fill=blue!8,align=center,font=\small}]
  \node[ct] (nft) {RoyaltyNFT\\[-2pt]\footnotesize ERC-721 + IMultiRoyalty};
  \node[ct,right=of nft] (mkt) {Marketplace\\[-2pt]\footnotesize Atomic Swap};
  \node[ct,below=of mkt] (spl) {PullPaymentSplitter\\[-2pt]\footnotesize Claim-based};
  \draw[-{Latex},thick] (mkt) -- node[above,font=\scriptsize]{royaltyShares} (nft);
  \draw[-{Latex},thick] (mkt) -- node[right,font=\scriptsize]{deposit} (spl);
  \draw[-{Latex},thick,dashed] (mkt.north east) .. controls +(1.5,0.8) and +(1.5,0.8) .. node[right,font=\scriptsize]{safeTransferFrom} (nft.north east);
\end{tikzpicture}

\vfill
\begin{tabular}{rl}
  \textbf{Stack:} & Solidity 0.8.24 $\cdot$ Hardhat 2.22 $\cdot$ OpenZeppelin 5 $\cdot$ ethers.js v6 \\
  \textbf{Chuẩn:} & ERC-721, EIP-2981, EIP-165 \\
  \textbf{Ngày:} & \today
\end{tabular}
\vspace{1cm}
\end{titlepage}

\tableofcontents
\newpage

% ==================== ABSTRACT ====================
\begin{abstract}
Báo cáo trình bày thiết kế, triển khai và phân tích một \emph{sàn giao dịch NFT với cơ chế chia
royalty phân mảnh} (\textit{Fractional Royalty NFT Marketplace}). Hệ thống gồm ba hợp đồng
thông minh: \textbf{(1)}~\texttt{RoyaltyNFT} — ERC-721 mở rộng EIP-2981 với giao diện
\texttt{IMultiRoyalty} cho phép nhiều stakeholder cùng nhận royalty;
\textbf{(2)}~\texttt{PullPaymentSplitter} — trình chia tiền theo mô hình pull-payment chống tấn
công tái nhập và cạn gas; \textbf{(3)}~\texttt{Marketplace} — thực hiện hoán đổi nguyên tử giữa
NFT và dòng tiền royalty trong một giao dịch duy nhất. Chúng tôi chứng minh toán học cho công thức
chia royalty theo basis points với xử lý dust, phân tích bảo mật theo SWC Registry, trình bày tối
ưu gas, và mô tả frontend tích hợp đầy đủ.
\end{abstract}
\newpage

% =====================================================================
\chapter{Giới thiệu}
\label{chap:intro}

\section{Bối cảnh}
Thị trường NFT đã trở thành cơ sở hạ tầng tài chính hoá tài sản số. Tuy nhiên, các marketplace
hiện nay gặp ba vấn đề cốt lõi:
\begin{enumerate}
  \item \textbf{Royalty một người nhận.} Chuẩn EIP-2981 chỉ hỗ trợ \emph{một} địa chỉ
        \texttt{receiver}, không đáp ứng nhu cầu chia sẻ giữa nhiều bên (nghệ sĩ, cộng tác viên,
        quỹ từ thiện, nhà đầu tư\ldots).
  \item \textbf{Rủi ro DoS khi chia tiền.} Cách tiếp cận ngây thơ dùng vòng lặp \texttt{.call}
        tới từng người nhận: một receiver xấu (\texttt{revert}, hết gas) phá vỡ mọi giao dịch.
  \item \textbf{Thiếu tính nguyên tử.} Nếu chuyển tiền và chuyển NFT tách rời, có thể xảy ra
        trạng thái trung gian: ``đã trả tiền nhưng chưa nhận NFT'' hoặc ngược lại.
\end{enumerate}

\section{Mục tiêu dự án}
Xây dựng hệ thống trong đó:
\begin{itemize}
  \item Mỗi NFT đăng ký danh sách \emph{nhiều} royalty receivers theo basis points (bps).
  \item Giao dịch bán diễn ra \emph{nguyên tử}: bất kỳ bước nào thất bại thì toàn bộ rollback.
  \item Tiền được chia theo \emph{pull-payment}: ghi credit trong mapping, người nhận tự withdraw.
\end{itemize}

\section{Phạm vi kỹ thuật}
\begin{itemize}
  \item Solidity 0.8.24, Hardhat 2.22, OpenZeppelin Contracts v5, evmVersion \texttt{cancun}.
  \item Frontend tĩnh: HTML + CSS + JavaScript thuần với ethers.js v6 (CDN, không build step).
  \item Mạng: Hardhat local (chainId 31337), cấu hình sẵn Sepolia testnet.
\end{itemize}

% =====================================================================
\chapter{Cơ sở lý thuyết}
\label{chap:background}

\section{ERC-721 và EIP-2981}
\textbf{ERC-721} định nghĩa NFT với hàm cốt lõi: \texttt{ownerOf}, \texttt{transferFrom},
\texttt{safeTransferFrom}, \texttt{approve}. Chúng tôi kế thừa \texttt{ERC721URIStorage} để lưu
metadata URI on-chain.

\textbf{EIP-2981} thêm duy nhất một hàm view:
\begin{lstlisting}[language=Solidity,numbers=none]
function royaltyInfo(uint256 tokenId, uint256 salePrice)
    external view returns (address receiver, uint256 royaltyAmount);
\end{lstlisting}
Hạn chế: chỉ một \texttt{receiver}. Để chia cho $n$ stakeholder, phải ủy thác qua contract
trung gian — tốn gas và ẩn thông tin.

\section{Giao diện \texttt{IMultiRoyalty} — Thiết kế mở rộng}

\begin{definition}[Multi-receiver royalty]\label{def:multiroy}
Với mỗi NFT có \texttt{tokenId} và giá bán $P$, hợp đồng trả về:
\begin{itemize}
  \item $T(P) = \lfloor P \cdot \alpha / 10000 \rfloor$: tổng royalty, với $\alpha \le 1000$ bps (tối đa 10\%).
  \item Mảng $\{(r_i, b_i)\}_{i=1}^{n}$ trong đó $r_i$ là địa chỉ, $b_i$ là phần chia (bps),
        thoả $\sum_{i=1}^{n} b_i = 10000$.
\end{itemize}
Mỗi người nhận $r_i$ được trả: $\text{amt}_i = \lfloor T \cdot b_i / 10000 \rfloor$.
\end{definition}

Struct \texttt{RoyaltyShare} sử dụng \texttt{uint96} cho \texttt{shareBps} để đóng gói đúng
256 bit cùng \texttt{address} (160 + 96 = 256), tiết kiệm storage slot.

\section{Pull-payment pattern}

\begin{definition}[Pull-payment]
Hợp đồng chỉ \emph{ghi công nợ} (credit) vào mapping $B: \texttt{address} \to \texttt{wei}$
khi nhận tiền; người nhận tự gọi \texttt{withdraw()} để rút.
\end{definition}

\begin{proposition}[Ưu thế chống DoS]\label{prop:pulldos}
Với push-payment (vòng lặp \texttt{.call} tới $n$ receiver), một receiver $r_k$ độc hại
(\texttt{fallback} revert hoặc tiêu hết gas) khiến toàn bộ giao dịch revert. Pull-payment cô
lập rủi ro: lỗi của $r_k$ chỉ ảnh hưởng đến $r_k$.
\end{proposition}

\section{CEI và ReentrancyGuard}

Tái nhập (\textit{reentrancy}) xảy ra khi external call trả quyền điều khiển cho kẻ tấn công
trong khi state chưa cập nhật. Dự án dùng \emph{hai lớp phòng thủ}:
\begin{enumerate}
  \item \textbf{CEI} (Checks-Effects-Interactions): sửa state trước external call.
  \item \textbf{ReentrancyGuard}: biến flag, nếu đã \texttt{ENTERED} thì revert.
\end{enumerate}

\section{Atomic swap trong EVM}

\begin{proposition}[Tính nguyên tử]\label{prop:atomic}
EVM đảm bảo: mọi revert rollback toàn bộ state của giao dịch. Do đó nếu chuyển tiền và chuyển
NFT nằm trong cùng một tx, bất kỳ bước nào revert sẽ hoàn tác tất cả — đạt tính nguyên tử
mà không cần cơ chế HTLC.
\end{proposition}

% =====================================================================
\chapter{Kiến trúc hệ thống}
\label{chap:arch}

\section{Tổng quan ba hợp đồng}

\begin{figure}[H]
\centering
\begin{tikzpicture}[
  every node/.style={font=\small},
  ct/.style={rectangle,rounded corners=5pt,draw,thick,minimum width=3.5cm,
             minimum height=1.3cm,fill=blue!5,align=center},
  ext/.style={rectangle,draw,thick,minimum width=3cm,minimum height=0.9cm,
              fill=orange!10,align=center,font=\small},
  arr/.style={-{Latex[length=2.5mm]},thick}
]
  \node[ct] (NFT) {RoyaltyNFT\\\footnotesize ERC-721 + ERC-2981\\\footnotesize + IMultiRoyalty};
  \node[ct,right=2.5cm of NFT] (MKT) {Marketplace\\\footnotesize ReentrancyGuard\\\footnotesize + Ownable};
  \node[ct,below=2cm of MKT] (SPL) {PullPaymentSplitter\\\footnotesize ReentrancyGuard};

  \node[ext,above=1.2cm of MKT] (FE) {Frontend\\[-2pt]\footnotesize HTML/JS + ethers.js v6};

  \draw[arr] (MKT) -- node[above,midway,font=\scriptsize]{royaltyShares()} (NFT);
  \draw[arr] (MKT) -- node[right,midway,font=\scriptsize]{deposit\{value\}} (SPL);
  \draw[arr,dashed] (MKT.south west) -- node[below,midway,sloped,font=\scriptsize]{safeTransferFrom} (NFT.south east);

  \draw[arr,blue!60] (FE) -- (MKT);
  \draw[arr,blue!60] (FE.west) -| node[left,near start,font=\scriptsize]{mint, approve} (NFT);
  \draw[arr,blue!60] (FE.east) -| node[right,near start,font=\scriptsize]{withdraw} (SPL);
\end{tikzpicture}
\caption{Kiến trúc hệ thống. Marketplace điều phối NFT và Splitter trong giao dịch atomic.
Frontend gọi trực tiếp cả ba hợp đồng.}
\label{fig:arch}
\end{figure}

\begin{table}[H]
\centering\small
\begin{tabular}{l l p{7.5cm}}
\toprule
\textbf{Contract} & \textbf{Kế thừa} & \textbf{Vai trò} \\
\midrule
\texttt{RoyaltyNFT} & ERC721URIStorage, ERC2981, Ownable & Mint NFT kèm mảng multi-receiver royalty \\
\texttt{Marketplace} & ReentrancyGuard, Ownable & List, buy (atomic swap), unlist \\
\texttt{PullPaymentSplitter} & ReentrancyGuard & Nhận deposit, ghi credit, cho phép withdraw \\
\bottomrule
\end{tabular}
\caption{Ba hợp đồng và vai trò tương ứng.}
\end{table}

\section{Sequence diagram: Atomic Buy}

\begin{figure}[H]
\centering
\begin{tikzpicture}[
  every node/.style={font=\scriptsize},
  actor/.style={rectangle,draw,minimum width=2cm,minimum height=0.7cm,fill=yellow!15},
  arr/.style={-{Latex},thick}
]
  \node[actor] (B) {Buyer};
  \node[actor,right=1.8cm of B] (M) {Marketplace};
  \node[actor,right=1.8cm of M] (N) {RoyaltyNFT};
  \node[actor,right=1.8cm of N] (S) {Splitter};

  \foreach \x in {B,M,N,S} {\draw[dashed] (\x) -- ++(0,-4);}

  \draw[arr] ($(B)+(0,-0.6)$) -- node[above]{\texttt{buy\{value:P\}}} ($(M)+(0,-0.6)$);
  \draw[arr] ($(M)+(0,-1.1)$) -- node[above]{\texttt{royaltyShares(id,P)}} ($(N)+(0,-1.1)$);
  \draw[arr] ($(N)+(0,-1.6)$) -- node[above]{(total, shares[])} ($(M)+(0,-1.6)$);
  \draw[arr] ($(M)+(0,-2.1)$) -- node[above]{\texttt{deposit\{value:P\}}} ($(S)+(0,-2.1)$);
  \draw[arr] ($(M)+(0,-2.6)$) -- node[above]{\texttt{safeTransferFrom}} ($(N)+(0,-2.6)$);
  \draw[arr] ($(N)+(0,-3.1)$) -- node[above]{NFT transferred} ($(B)+(0,-3.1)$);

  \draw[thick,decorate,decoration={brace,mirror,amplitude=5pt}]
    ($(B)+(-0.5,-0.4)$) -- ($(B)+(-0.5,-3.4)$) node[midway,left=8pt,font=\scriptsize,align=right]{Một TX\\duy nhất};
\end{tikzpicture}
\caption{Sequence diagram \texttt{Marketplace.buy()}: tất cả nằm trong một giao dịch duy nhất.}
\label{fig:seq}
\end{figure}

\section{Luồng hoạt động tổng quát}
\begin{enumerate}
  \item \textbf{Mint}: Owner gọi \texttt{mintWithRoyalties(to, uri, royaltyBps, shares[])}.
  \item \textbf{Approve}: Seller gọi \texttt{nft.approve(marketplace, tokenId)}.
  \item \textbf{List}: Seller gọi \texttt{marketplace.list(nft, tokenId, price)}.
  \item \textbf{Buy (atomic)}: Buyer gọi \texttt{marketplace.buy(nft, tokenId)} với \texttt{msg.value = price}:
    \begin{itemize}
      \item Xoá listing (Effects trước Interactions).
      \item Tính fee + royalty + seller proceeds.
      \item \texttt{splitter.deposit\{value: price\}(payees, amounts)} — ghi credit.
      \item \texttt{nft.safeTransferFrom(seller, buyer, tokenId)} — chuyển NFT.
      \item Nếu bất kỳ bước nào revert $\Rightarrow$ toàn bộ rollback.
    \end{itemize}
  \item \textbf{Withdraw}: Mỗi stakeholder tự gọi \texttt{splitter.withdraw()} để rút ETH.
\end{enumerate}

% =====================================================================
\chapter{Cơ sở toán học và chứng minh}
\label{chap:math}

\section{Công thức chia royalty theo basis points}

\begin{definition}[Hệ thống basis points]
$1\text{ bps} = 0.01\%$. Denominator $D = 10000$. Royalty rate $\alpha$ bps trên giá $P$:
\[
  T = \left\lfloor \frac{P \cdot \alpha}{D} \right\rfloor, \qquad \alpha \le 1000 \text{ (tối đa 10\%)}.
\]
\end{definition}

\begin{definition}[Phân phối cho receivers]
Với mảng shares $(b_1, \ldots, b_n)$ thoả $\sum_{i=1}^{n} b_i = D = 10000$:
\[
  \text{amt}_i = \left\lfloor \frac{T \cdot b_i}{D} \right\rfloor, \qquad i = 1, \ldots, n.
\]
\end{definition}

\section{Xử lý sai số làm tròn (dust)}

\begin{theorem}[Bảo toàn tổng với dust]\label{thm:dust}
Gọi $d = T - \sum_{i=1}^{n} \lfloor T \cdot b_i / D \rfloor$ là phần dư (dust).
Khi đó $0 \le d < n$ và thuật toán cộng $d$ vào receiver cuối đảm bảo $\sum \text{amt}_i = T$.
\end{theorem}
\begin{proof}
Đặt $a_i = \lfloor T b_i / D \rfloor$. Ta có $a_i \le T b_i / D$, nên:
\[
  \sum_{i=1}^{n} a_i \le \sum_{i=1}^{n} \frac{T b_i}{D} = \frac{T}{D}\sum b_i = T.
\]
Mặt khác, $T b_i / D - a_i < 1$ nên $d = T - \sum a_i < n$. Trong mã nguồn:
\begin{lstlisting}[language=Solidity,numbers=none]
if (n > 0 && distributed < total) {
    amounts[n - 1] += (total - distributed); // dust -> last receiver
}
\end{lstlisting}
Sau bước này, $\sum \text{amt}_i = T$ chính xác.
\end{proof}

\section{Phân phối tổng giá bán}

\begin{proposition}[Bảo toàn giá bán]\label{prop:priceconserve}
Với giá $P$, phí marketplace $f$ bps, và royalty $\alpha$ bps:
\[
  P = \underbrace{\lfloor P \cdot f / D \rfloor}_{\text{fee}} +
      \underbrace{\lfloor P \cdot \alpha / D \rfloor}_{\text{royalty}} +
      \text{proceeds}_{\text{seller}}.
\]
Trong mã nguồn, \texttt{proceeds = price - fee - totalRoyalty} tính trực tiếp phần dư, đảm bảo
không mất wei nào.
\end{proposition}

\section{Tính đúng đắn của deposit}

\begin{lemma}[Bất biến PullPaymentSplitter]\label{lem:splitter-inv}
Tại mọi thời điểm:
\[
  \texttt{address(this).balance} \ge \sum_{\text{addr}} \texttt{\_balances[addr]} = \texttt{\_totalPending}.
\]
\end{lemma}
\begin{proof}
Hàm \texttt{deposit} kiểm tra $\sum \texttt{amounts} = \texttt{msg.value}$, nên mỗi lần deposit,
số dư ETH tăng đúng bằng tổng credit thêm vào mapping. Hàm \texttt{withdraw} trừ balance trước
khi gửi ETH (CEI), nên hai vế giảm đồng bộ.
\end{proof}

% =====================================================================
\chapter{Phân tích mã nguồn chi tiết}
\label{chap:code}

\section{\texttt{IMultiRoyalty.sol} — Giao diện mở rộng EIP-2981}

\begin{lstlisting}[language=Solidity,caption={Giao diện IMultiRoyalty.}]
interface IMultiRoyalty {
    struct RoyaltyShare {
        address receiver;  // 160 bits
        uint96 shareBps;   //  96 bits -> total 256 bits = 1 slot
    }
    function royaltyShares(uint256 tokenId, uint256 salePrice)
        external view
        returns (uint256 totalAmount, RoyaltyShare[] memory shares);
}
\end{lstlisting}

\textbf{Thiết kế}: \texttt{address} (160 bit) + \texttt{uint96} (96 bit) = 256 bit, đóng gói vừa
một storage slot. Marketplace chỉ cần gọi \emph{một hàm} để lấy cả tổng royalty lẫn danh sách chia.
Đây là điểm khác biệt so với EIP-2981 chuẩn (chỉ 1 receiver).

\section{\texttt{RoyaltyNFT.sol} — ERC-721 với multi-receiver royalty}

\begin{lstlisting}[language=Solidity,caption={Hàm \texttt{mintWithRoyalties} — mint kèm cấu hình royalty.}]
function mintWithRoyalties(
    address to, string calldata tokenURI_,
    uint96 royaltyBps, RoyaltyShare[] calldata shares
) external onlyOwner returns (uint256 tokenId) {
    if (royaltyBps > MAX_ROYALTY_BPS) revert RoyaltyTooHigh(royaltyBps);
    if (shares.length == 0) revert EmptyShares();
    tokenId = ++_nextTokenId;
    _safeMint(to, tokenId);
    _setTokenURI(tokenId, tokenURI_);
    _setRoyaltyShares(tokenId, royaltyBps, shares);
    _setTokenRoyalty(tokenId, address(this), royaltyBps);
}
\end{lstlisting}

\textbf{Giải thích}:
\begin{itemize}
  \item \texttt{MAX\_ROYALTY\_BPS = 1000} (10\%): giới hạn cứng ngăn owner set royalty quá cao.
  \item \texttt{\_setTokenRoyalty(tokenId, address(this), royaltyBps)}: đặt \texttt{receiver = address(this)}
        trong EIP-2981 chuẩn — tương thích ngược với marketplace bên ngoài.
  \item \texttt{\_setRoyaltyShares}: lưu mảng shares vào mapping, kiểm tra $\sum b_i = 10000$
        và $r_i \ne \texttt{address(0)}$.
  \item \texttt{supportsInterface} override trả \texttt{true} cho cả ERC-2981 và \texttt{IMultiRoyalty}
        — giúp Marketplace dò qua EIP-165 trước khi gọi.
\end{itemize}

\section{\texttt{PullPaymentSplitter.sol} — Claim-based payment}

\begin{lstlisting}[language=Solidity,caption={Deposit (chỉ ghi mapping) và Withdraw (CEI).}]
function deposit(address[] calldata payees, uint256[] calldata amounts)
    external payable
{
    require(payees.length == amounts.length, "length mismatch");
    uint256 total;
    for (uint256 i = 0; i < payees.length; ++i) {
        if (payees[i] == address(0)) revert ZeroPayee();
        _balances[payees[i]] += amounts[i];
        total += amounts[i];
        emit PaymentDeposited(payees[i], amounts[i], msg.sender);
    }
    if (total != msg.value) revert AmountMismatch(total, msg.value);
    _totalPending += total;
}

function _withdrawTo(address payable payee) internal {
    uint256 amount = _balances[payee];
    if (amount == 0) revert NothingToWithdraw();
    _balances[payee] = 0;          // Effects TRUOC Interactions
    _totalPending -= amount;
    emit PaymentWithdrawn(payee, amount);
    Address.sendValue(payee, amount); // Interactions
}
\end{lstlisting}

\textbf{Ba điểm bảo mật quan trọng}:
\begin{enumerate}
  \item \texttt{deposit} chỉ ghi mapping — \emph{không external call} cho từng payee
        $\Rightarrow$ không DoS, không reentrancy.
  \item \texttt{\_withdrawTo} reset balance = 0 \emph{trước} \texttt{sendValue} (CEI).
        Nếu bị tái nhập, lần gọi thứ hai revert tại \texttt{NothingToWithdraw}.
  \item \texttt{Address.sendValue} dùng \texttt{call\{value\}("")} (toàn bộ gas) — tương thích
        smart-contract wallets (Gnosis Safe) vốn tốn $>$2300 gas.
\end{enumerate}

\section{\texttt{Marketplace.sol} — Atomic swap}

\begin{lstlisting}[language=Solidity,caption={Hàm \texttt{buy()} — trái tim của atomic swap.}]
function buy(address nft, uint256 tokenId)
    external payable nonReentrant
{
    Listing memory l = _listings[nft][tokenId];
    if (!l.active)            revert NotListed();     // Check
    if (msg.value != l.price) revert IncorrectPayment(l.price, msg.value);
    if (msg.sender == l.seller) revert CannotBuyOwn();

    delete _listings[nft][tokenId]; // Effect: xoa TRUOC interactions

    uint256 price = l.price;
    uint256 fee = (price * marketplaceFeeBps) / BPS;
    (uint256 totalRoyalty, address[] memory rcv, uint256[] memory amt) =
        _computeRoyalty(nft, tokenId, price);
    uint256 proceeds = price - fee - totalRoyalty;

    // Gom tat ca payees thanh 1 mang duy nhat
    uint256 n = rcv.length;
    address[] memory payees = new address[](n + 2);
    uint256[] memory amounts = new uint256[](n + 2);
    for (uint256 i = 0; i < n; ++i) {
        payees[i] = rcv[i]; amounts[i] = amt[i];
    }
    payees[n] = feeRecipient;  amounts[n] = fee;
    payees[n+1] = l.seller;    amounts[n+1] = proceeds;

    splitter.deposit{value: price}(payees, amounts); // Interaction (a)
    IERC721(nft).safeTransferFrom(                   // Interaction (b)
        l.seller, msg.sender, tokenId);

    emit Sold(nft, tokenId, msg.sender, l.seller,
              price, totalRoyalty, fee, proceeds);
}
\end{lstlisting}

\textbf{Royalty resolver} (\texttt{\_computeRoyalty}): dùng \texttt{try/catch} hai tầng:
\begin{enumerate}
  \item Thử \texttt{IMultiRoyalty} qua EIP-165 \texttt{supportsInterface} — nếu OK, gọi
        \texttt{royaltyShares()} để lấy multi-receiver.
  \item Fallback \texttt{IERC2981.royaltyInfo()} — 1 receiver.
  \item Nếu cả hai fail (NFT không hỗ trợ royalty): royalty = 0.
\end{enumerate}
Thiết kế \texttt{try/catch} đảm bảo NFT độc hại (revert ngẫu nhiên) không phá vỡ Marketplace.

% =====================================================================
\chapter{Phân tích bảo mật}
\label{chap:security}

\section{Ma trận đe doạ theo SWC Registry}

\begin{table}[H]
\centering\small
\begin{tabularx}{\textwidth}{p{3.8cm} l X}
\toprule
\textbf{Mối đe doạ} & \textbf{SWC} & \textbf{Biện pháp} \\
\midrule
Reentrancy \texttt{buy()} & 107 & \texttt{nonReentrant} + CEI (xoá listing trước external call) \\
Reentrancy \texttt{withdraw()} & 107 & \texttt{nonReentrant} + reset balance = 0 trước \texttt{sendValue} \\
DoS chia tiền nhiều người & 113 & Pull-payment: deposit chỉ ghi mapping, không external call \\
Out-of-gas fallback & 126 & Cô lập trong withdraw riêng từng receiver \\
Tràn số học & 101 & Solidity 0.8+ auto-check overflow/underflow \\
Front-running unlist/relist & 114 & \texttt{IncorrectPayment}: buyer phải trả đúng giá listed \\
NFT contract độc hại & --- & \texttt{try/catch} trong \texttt{\_computeRoyalty} \\
Owner rug-pull phí & --- & Giới hạn cứng: \texttt{marketplaceFeeBps} $\le$ 1000 (10\%) \\
\bottomrule
\end{tabularx}
\caption{Ma trận đe doạ và biện pháp phòng vệ.}
\end{table}

\section{Chứng minh an toàn reentrancy của \texttt{buy()}}

\begin{proposition}
Hàm \texttt{Marketplace.buy()} không thể bị khai thác qua tái nhập.
\end{proposition}
\begin{proof}
Hai điểm external call: (a)~\texttt{splitter.deposit}, (b)~\texttt{safeTransferFrom}.
\begin{itemize}
  \item \textit{Trước (a)}: listing đã bị \texttt{delete}. Re-entry vào \texttt{buy(nft, tokenId)}
        revert tại \texttt{NotListed}.
  \item \textit{Tại (a)}: \texttt{deposit} chỉ ghi mapping, không callback.
  \item \textit{Tại (b)}: \texttt{safeTransferFrom} gọi \texttt{onERC721Received} của buyer.
        Nếu buyer là contract độc tái nhập \texttt{buy(nft', tokenId')} khác,
        \texttt{nonReentrant} là guard \emph{toàn cục} cho hàm $\Rightarrow$ revert.
\end{itemize}
Kết luận: không có trạng thái trung gian nào exploitable.
\end{proof}

\section{Chứng minh an toàn pull-payment}

\begin{proposition}
Không thể rút quá balance đã được ghi nhận.
\end{proposition}
\begin{proof}
\texttt{\_withdrawTo(payee)}: đọc $a = \texttt{\_balances[payee]}$, gán
$\texttt{\_balances[payee]} = 0$ (Effects), rồi \texttt{sendValue(payee, a)} (Interactions).
Nếu \texttt{sendValue} bị tái nhập, vòng thứ hai đọc $a = 0$ $\Rightarrow$
\texttt{NothingToWithdraw} revert. Tổng ETH rút $\le$ tổng deposit vì balances chỉ tăng qua
\texttt{deposit}, chỉ giảm qua \texttt{withdraw}.
\end{proof}

% =====================================================================
\chapter{Tối ưu Gas}
\label{chap:gas}

\section{Kỹ thuật áp dụng}

\begin{table}[H]
\centering\small
\begin{tabularx}{\textwidth}{l X}
\toprule
\textbf{Kỹ thuật} & \textbf{Chi tiết} \\
\midrule
Storage packing & \texttt{RoyaltyShare}: \texttt{address} (160) + \texttt{uint96} (96) = 1 slot \\
Custom errors & Thay \texttt{revert("string")} bằng \texttt{revert ErrorName()} — tiết kiệm $\sim$50 gas/call và bytecode deploy \\
Immutable & \texttt{splitter} trong Marketplace là \texttt{immutable} — đọc không SLOAD \\
Gom deposit & Một lần \texttt{splitter.deposit} cho tất cả payees thay vì nhiều lần \texttt{call} \\
unchecked \texttt{++i} & Biến vòng lặp tăng trong \texttt{unchecked} (đã kiểm tra bound) \\
via-IR + optimizer & \texttt{hardhat.config.js}: \texttt{viaIR: true}, \texttt{runs: 200} \\
\bottomrule
\end{tabularx}
\caption{Các kỹ thuật tối ưu gas.}
\end{table}

\section{Phân tích gas đo được}

Số liệu từ test suite trên Hardhat local:
\begin{table}[H]
\centering\small
\begin{tabular}{l l r}
\toprule
\textbf{Contract} & \textbf{Method} & \textbf{Gas (approx)} \\
\midrule
RoyaltyNFT & \texttt{mintWithRoyalties} (3 receivers) & $\sim$264\,000 \\
RoyaltyNFT & \texttt{approve} & $\sim$48\,000 \\
Marketplace & \texttt{list} & $\sim$60\,000 \\
Marketplace & \texttt{buy} (atomic, 3 royalty receivers) & $\sim$245\,000 \\
PullPaymentSplitter & \texttt{deposit} (5 payees) & $\sim$94\,000 \\
PullPaymentSplitter & \texttt{withdraw} & $\sim$37\,000 \\
\bottomrule
\end{tabular}
\caption{Gas consumption các hàm chính.}
\end{table}

\textbf{Nhận xét}: \texttt{buy()} tốn $\sim$245k gas (bao gồm cả deposit + safeTransferFrom) — nằm
thoải mái trong block gas limit (30M), cho phép nhiều giao dịch cùng block.

% =====================================================================
\chapter{Kiểm thử}
\label{chap:testing}

\section{Chiến lược}

Bộ test sử dụng Hardhat + Chai + ethers.js, bao gồm:
\begin{enumerate}
  \item \textbf{Unit test}: kiểm tra từng hàm riêng lẻ.
  \item \textbf{Integration test}: luồng đầy đủ mint $\to$ approve $\to$ list $\to$ buy $\to$ withdraw.
  \item \textbf{Negative test}: kiểm tra revert khi input không hợp lệ.
\end{enumerate}

\section{Các test case}

\begin{table}[H]
\centering\small
\begin{tabularx}{\textwidth}{c p{3cm} X}
\toprule
\textbf{\#} & \textbf{Test} & \textbf{Ý nghĩa} \\
\midrule
1 & Mint multi-receiver & Mint NFT với 3 receivers, kiểm tra \texttt{royaltyShares()} trả đúng tổng 10\% và 3 shares \\
2 & Atomic sale & Approve $\to$ list $\to$ buy: NFT chuyển cho buyer, fee/royalty/proceeds ghi đúng credit trong splitter \\
3 & Pull-payment & Mỗi stakeholder tự withdraw, kiểm tra ETH balance tăng đúng số expected \\
4 & Incorrect payment & Gửi sai giá $\Rightarrow$ revert \texttt{IncorrectPayment} \\
\bottomrule
\end{tabularx}
\caption{Bốn test case chính.}
\end{table}

\section{Ví dụ test: Atomic sale với chia royalty}

\begin{lstlisting}[language=JavaScript,caption={Test kiểm tra chia tiền đúng sau atomic buy.}]
it("executes atomic sale: NFT moves, funds credited", async () => {
  const price = ethers.parseEther("1");
  // Seller approve + list
  await nft.connect(seller).approve(market, 1);
  await market.connect(seller).list(nft, 1, price);
  // Buyer buys
  await expect(market.connect(buyer).buy(nft, 1, { value: price }))
    .to.emit(market, "Sold");
  // NFT transferred
  expect(await nft.ownerOf(1)).to.equal(buyer.address);
  // Fee 2.5% = 0.025 ETH, Royalty 10% = 0.1 ETH
  // Artist 50%: 0.05, Charity 30%: 0.03, Investor 20%: 0.02
  // Seller proceeds: 1 - 0.025 - 0.1 = 0.875 ETH
  expect(await splitter.balanceOf(artist.address))
    .to.equal(ethers.parseEther("0.05"));
  expect(await splitter.balanceOf(charity.address))
    .to.equal(ethers.parseEther("0.03"));
  expect(await splitter.balanceOf(seller.address))
    .to.equal(ethers.parseEther("0.875"));
});
\end{lstlisting}

\section{Kết quả}

\begin{verbatim}
  NFT Marketplace with Fractional Royalties
    [PASS] mints NFT with multi-receiver royalty metadata (734ms)
    [PASS] executes atomic sale: NFT moves, funds credited (73ms)
    [PASS] allows each stakeholder to pull-payment independently (48ms)
    [PASS] reverts atomically when payment is incorrect (62ms)

  4 passing (927ms)
\end{verbatim}

% =====================================================================
\chapter{Frontend}
\label{chap:frontend}

\section{Kiến trúc}

Frontend là ứng dụng tĩnh thuần HTML/CSS/JavaScript, không framework, không build step.
Dependencies duy nhất: ethers.js v6 qua CDN.

\begin{table}[H]
\centering\small
\begin{tabular}{l p{9cm}}
\toprule
\textbf{File} & \textbf{Vai trò} \\
\midrule
\texttt{index.html} & Cấu trúc giao diện: 4 tab (Marketplace, My NFTs, Mint, Withdraw) \\
\texttt{index.css} & Design system: dark-mode, glassmorphism, gradients, responsive \\
\texttt{app.js} & Toàn bộ logic: wallet connect, contract interaction, event listeners \\
\bottomrule
\end{tabular}
\end{table}

\section{Tích hợp Smart Contract}

ABI dạng Human-Readable của ethers v6 giúp code gọn:
\begin{lstlisting}[language=JavaScript,numbers=none]
const MARKETPLACE_ABI = [
  "function buy(address nft, uint256 tokenId) payable",
  "function list(address nft, uint256 tokenId, uint256 price)",
  "function getListing(address,uint256) view returns (tuple(address,uint256,bool))",
  "event Sold(address indexed nft, uint256 indexed tokenId, ...)"
];
const contract = new ethers.Contract(addr, MARKETPLACE_ABI, signer);
\end{lstlisting}

\section{Tự động phát hiện contracts}

Thay vì hardcode địa chỉ, \texttt{app.js} quét blockchain để tìm contracts đã deploy:
\begin{enumerate}
  \item Scan các block gần đây tìm contract creation transactions từ deployer.
  \item Thử gọi \texttt{name()} (RoyaltyNFT) và \texttt{splitter()} (Marketplace) để xác định loại.
  \item Nếu không tìm được, hiện prompt nhập thủ công.
\end{enumerate}

\section{Event-driven UI}

Frontend đăng ký listen contract events:
\begin{itemize}
  \item \texttt{Transfer(from=0x0, ...)} $\Rightarrow$ Toast ``NFT Minted!''
  \item \texttt{Listed(...)} $\Rightarrow$ Refresh marketplace grid
  \item \texttt{Sold(...)} $\Rightarrow$ Toast + refresh cả marketplace và balance
\end{itemize}

% =====================================================================
\chapter{Hướng dẫn chạy và demo}
\label{chap:run}

\section{Cài đặt}
\begin{lstlisting}[numbers=none]
git clone <repo-url>
cd smart-contract
npm install
npx hardhat compile
npx hardhat test          # chay 4 tests
\end{lstlisting}

\section{Demo trên local}
Mở ba terminal:
\begin{lstlisting}[numbers=none]
# Terminal 1: Hardhat node (20 accounts x 10000 ETH)
npx hardhat node

# Terminal 2: Deploy + seed demo data
npm run deploy:local
npm run seed:local        # mint 3 NFTs, list 2

# Terminal 3: Frontend
npm run serve             # http://localhost:3000
\end{lstlisting}

\section{Kịch bản demo 5 phút}
\begin{enumerate}
  \item \textbf{Connect}: Mở trình duyệt $\to$ Click ``Connect Wallet'' $\to$ auto-detect contracts.
  \item \textbf{Mint}: Tab Mint $\to$ nhập URI, royalty 5\%, 2 receivers (70/30) $\to$ Mint.
  \item \textbf{List}: Tab My NFTs $\to$ click NFT card $\to$ ``List for Sale'' $\to$ nhập 1 ETH.
  \item \textbf{Buy}: Đổi account MetaMask $\to$ Tab Marketplace $\to$ click NFT $\to$ ``Buy''.
  \item \textbf{Verify}: Kiểm tra NFT đã chuyển owner. Tab Withdraw $\to$ thấy balance pending.
  \item \textbf{Withdraw}: Click ``Withdraw All'' $\to$ ETH về ví.
\end{enumerate}

% =====================================================================
\chapter{Kết luận}
\label{chap:conclusion}

\section{Kết quả đạt được}

Dự án đã triển khai đầy đủ ba yêu cầu kỹ thuật:

\begin{table}[H]
\centering\small
\begin{tabularx}{\textwidth}{p{4cm} l X}
\toprule
\textbf{Yêu cầu} & \textbf{Contract} & \textbf{Đặc điểm nổi bật} \\
\midrule
Advanced ERC-721 + EIP-2981 multi-receiver & RoyaltyNFT & Giao diện \texttt{IMultiRoyalty} chuẩn hoá; storage-packed struct; tương thích ngược EIP-2981 \\
Pull-payment chống Reentrancy \& DoS & PullPaymentSplitter & CEI + ReentrancyGuard; deposit không external call; \texttt{Address.sendValue} tương thích contract wallets \\
Atomic swap & Marketplace & Một tx duy nhất: deposit + safeTransferFrom; \texttt{try/catch} cho NFT độc hại; custom errors \\
\bottomrule
\end{tabularx}
\end{table}

\section{Điểm mạnh kỹ thuật}
\begin{itemize}
  \item \textbf{Custom logic}: Giao diện \texttt{IMultiRoyalty} là thiết kế riêng, mở rộng EIP-2981
        theo hướng chưa chuẩn hoá — không thể clone từ boilerplate OpenZeppelin.
  \item \textbf{Royalty resolver 2 tầng}: \texttt{try/catch} cho cả \texttt{IMultiRoyalty} lẫn
        EIP-2981 chuẩn, fallback gracefully khi NFT contract không hỗ trợ.
  \item \textbf{Bảo mật nhiều lớp}: CEI + ReentrancyGuard + pull-payment + custom errors +
        giới hạn cứng (max royalty, max fee).
  \item \textbf{Frontend thuần}: không dependency nặng, không build step — dễ deploy, dễ audit.
\end{itemize}

\section{Hướng phát triển}
\begin{enumerate}
  \item Hỗ trợ auction (English/Dutch).
  \item Thanh toán bằng ERC-20.
  \item Batch buy/list nhiều NFT trong 1 tx.
  \item Off-chain signed orders (kiểu Seaport).
  \item Tích hợp IPFS/Arweave cho metadata.
\end{enumerate}

% =====================================================================
\begin{thebibliography}{9}
\bibitem{eip721}
  W.~Entriken et al.,
  ``EIP-721: Non-Fungible Token Standard,'' 2018.
  \url{https://eips.ethereum.org/EIPS/eip-721}

\bibitem{eip2981}
  Z.~Goldberg, J.~Morgan,
  ``EIP-2981: NFT Royalty Standard,'' 2020.
  \url{https://eips.ethereum.org/EIPS/eip-2981}

\bibitem{eip165}
  C.~Burgess et al.,
  ``EIP-165: Standard Interface Detection,'' 2018.
  \url{https://eips.ethereum.org/EIPS/eip-165}

\bibitem{oz}
  OpenZeppelin,
  ``OpenZeppelin Contracts v5.''
  \url{https://docs.openzeppelin.com/contracts/5.x/}

\bibitem{swc}
  SmartContractSecurity,
  ``SWC Registry.''
  \url{https://swcregistry.io/}

\bibitem{pullpayment}
  ConsenSys,
  ``Ethereum Smart Contract Best Practices: Pull Payment.''
  \url{https://consensys.github.io/smart-contract-best-practices/}
\end{thebibliography}

\end{document}
